The Convex Hull Ratio is a uniquely valuable compactness metric for redistricting
because it detects non-convex appendages that perimeter-based and global area-based
metrics fail to identify.
Our principal findings are:

\begin{enumerate}
  \item Synthetic tentacle districts with moderate PP ($\approx 0.27$) and
    moderate Reock ($\approx 0.18$) can score CH $\approx 0.52$, correctly
    flagging the non-convex geometry that PP and Reock underweight.
  \item All four bisect structure algorithms guarantee CH $> 0.85$ on all
    NC 2020 congressional districts---a structural property of recursive
    graph bisection that ensures no tentacle districts are produced.
  \item CH provides the most legally intuitive courtroom presentation:
    the ``rubber band'' explanation requires no mathematical training,
    and the convex hull overlay directly operationalises the visual-irregularity
    test in \textit{Shaw v.\ Reno} (1993).
\end{enumerate}

Practitioners should include CH in all compactness profiles, particularly
when the legal challenge involves a claim of ``bizarre'' or ``irregular'' shape
under state constitutional or statutory compactness requirements.
The K.7 composite guide integrates CH into a certified six-metric profile.
